\section{Limitações e Evolução}
\label{sec:limitacoes}

\subsection{O que o validador não faz}

O validador responde a uma pergunta específica: ``esta sequência de
caracteres tem o formato de uma placa brasileira?''. Ele não
responde a outras perguntas que podem parecer próximas.

\begin{table}[ht]
\centering
\small
\renewcommand{\arraystretch}{1.3}
\begin{tabular}{p{4.2cm} p{4.2cm} p{4.6cm}}
\toprule
Limitação & Consequência & Como tratar \\
\midrule
Não consulta cadastros
  & aceita \texttt{ABC1234} mesmo que nenhum veículo tenha essa placa
  & consulta ao órgão de trânsito, depois da validação do formato \\
Não distingue placa convertida de placa nova
  & aceita qualquer letra de A a Z na 5ª posição do Mercosul
  & usar o AFD 03, como verificação adicional \\
AFDs puros aceitam só maiúsculas e sem espaços internos
  & \texttt{abc1234} (nos AFDs) e \texttt{ABC 1234} são rejeitadas
  & manter pré-normalização (\texttt{upper/strip}) e remover espaços internos \\
Cobre apenas dois formatos
  & placas fora desses formatos são rejeitadas
  & acrescentar um autômato por formato novo \\
Autômato produto mínimo dá resposta binária
  & o produto de 11 estados não informa qual padrão foi reconhecido nem onde está o erro
  & executar AFDs separados (como em \texttt{identificar\_modelo}) ou emitir mensagem por estado \\
\bottomrule
\end{tabular}
\caption{Limitações do validador e formas de tratá-las.}
\label{tab:limitacoes}
\end{table}

\subsection{Existência da placa}

Reconhecer o formato não prova que a placa exista. A quantidade de
cadeias aceitas ($808\,496\,000$, contando as duas grafias do
padrão anterior) é muito maior que a de placas emitidas, e saber
quais foram emitidas depende de bases de dados que um autômato não
tem como conter. Por isso o validador deve ser usado como
\emph{primeira etapa}: ele barra erros de digitação com resposta
imediata, e só as placas aprovadas seguem para a consulta ao
cadastro.

\subsection{Decisões que dependem da regra vigente}

Três decisões do projeto refletem a interpretação adotada neste
trabalho, e podem mudar se a regra ou o uso mudarem:

\begin{itemize}
  \item o hífen é aceito de forma opcional no padrão anterior,
        embora ele seja apenas uma convenção visual da chapa;
  \item a quinta posição do Mercosul aceita qualquer letra, pois o
        validador não sabe se a placa foi convertida;
  \item os autômatos puros operam apenas sobre letras em caixa alta,
        sem acento, delegando a normalização à camada de entrada.
\end{itemize}

As normas que definem as placas são resoluções do CONTRAN
\cite{contran780,contran741} e podem ser alteradas. Se isso
acontecer, basta ajustar a especificação, refazer a tabela de
transição e repetir a verificação. O projeto foi organizado para
que essa revisão seja pequena.

\subsection{Limite do modelo}

Todo o conjunto de placas é finito, então qualquer linguagem desse
tipo é regular e o autômato finito é suficiente. Um modelo mais
poderoso só seria necessário para regras que exigem memória
ilimitada, o que não é o caso aqui. Também por isso o autômato
pequeno (11 estados no produto mínimo) consegue substituir uma lista
com mais de seiscentos milhões de placas possíveis.

\subsection{Evolução possível}

\begin{itemize}
  \item \textbf{Normalização completa da entrada:} complementar a
        conversão para maiúsculas e remoção de bordas já existente
        em \texttt{identificar\_modelo} com a remoção de espaços
        internos (como em \texttt{ABC 1234}).
  \item \textbf{Classificação em autômato único:} utilizar a versão
        de 14 estados do autômato produto (com estados finais
        separados) para identificar o padrão em uma única passada,
        sem precisar executar dois autômatos sequencialmente.
  \item \textbf{Validação durante a digitação:} como cada estado
        representa uma etapa da placa, o validador pode avisar a
        cada caractere se o que foi digitado ainda pode formar uma
        placa válida.
  \item \textbf{Mensagens de erro específicas:} o estado em que a
        leitura falha indica o que estava esperado, por exemplo
        ``falta um algarismo na 4ª posição''.
  \item \textbf{Extração em textos livres:} procurar placas dentro
        de um texto, como o de uma multa ou de um e-mail, usando o
        reconhecimento do maior trecho possível.
  \item \textbf{Integração com cadastros:} encadear o validador com
        a consulta ao órgão de trânsito.
\end{itemize}